\lysh{16153}{4}

\begin{alist}
\item Πρέπει $x \ne 0$, άρα $A_f = \mathbb{R}^* = (-\infty, 0) \cup (0, +\infty)$.

\item Αφού είναι $x \ne 0$, η γραφική παράσταση δεν μπορεί να τέμνει τον άξονα $y'y$ αφού τα
σημεία του $y'y$ έχουν τετμημένη μηδέν.
Αν υποθέσουμε ότι η γραφική παράσταση της $f$ τέμνει τον άξονα $x'x$ σε σημείο με τετμημένη
$\beta$, τότε θα έπρεπε $f(\beta) = 0$, δηλαδή $1 + \dfrac{4}{\beta^2} = 0 \Leftrightarrow \dfrac{4}{\beta^2} = -1 \Leftrightarrow \beta^2 = -4$, αδύνατο, άρα η
γραφική παράσταση της $f$ δεν τέμνει ούτε τον άξονα $x'x$.

\item 
\begin{rlist}
\item Το σημείο $M$ έχει συντεταγμένες $M(\alpha, f(\alpha))$, οπότε το εμβαδόν του ορθογωνίου θα είναι
\[E = (O\Delta) \cdot (M\Delta) = \alpha \cdot f(\alpha) = \alpha \left(1 + \dfrac{4}{\alpha^2}\right) = \alpha \cdot 1 + \alpha \cdot \dfrac{4}{\alpha^2} = \alpha + \dfrac{4}{\alpha}.\]

\item $E \ge 4 \Leftrightarrow \alpha + \dfrac{4}{\alpha} \ge 4 \Leftrightarrow \alpha^2 + 4 \ge 4\alpha \Leftrightarrow \alpha^2 + 4 - 4\alpha \ge 0 \Leftrightarrow (\alpha - 2)^2 \ge 0$, ισχύει
για κάθε πραγματικό αριθμό $\alpha$, οπότε και για $\alpha > 0$.
\end{rlist}
\end{alist}

